\section{Thảo luận và giới hạn}

Kết quả Q1--Q3 nhất quán với nhau: chuỗi có xu hướng dài hạn mạnh, mùa vụ yếu và các
đợt cực trị tập trung ở một số năm. Chính vì mùa vụ yếu, Seasonal Naive 365 thất bại
nặng; và vì chuỗi gần random-walk ở tần suất ngày, Naive persistence là baseline mạnh.

Các giới hạn chính:
\begin{itemize}
  \item Chỉ có 14 năm dữ liệu ngày; Test chỉ 2 năm, chưa đủ cho kết luận dài hạn.
  \item Chưa dùng biến ngoại sinh như mưa, nhiệt độ, tuyết tan hay vận hành hồ chứa.
  \item GRU chưa chạy trong môi trường hiện tại; báo cáo không có kết quả GRU.
  \item LightGBM dùng mục tiêu delta; đây là lựa chọn kỹ thuật đã ghi rõ, không phải
        kết quả được chọn theo Test.
  \item Đây là backtest, không phải dự báo vận hành hoặc cảnh báo lũ.
\end{itemize}

Về mặt phương pháp, việc giữ Test tách biệt và không tinh chỉnh trên Test là điều kiện
cần để kết quả có ý nghĩa. Việc Naive thắng LightGBM một chút là kết quả trung thực,
không phải lỗi cần che giấu. Nếu tiếp tục phát triển, nên thử mô hình có biến ngoại
sinh, kiểm định rolling-origin dài hơn và phân tích sai số theo mùa.
